%==============================================================================
% CAPÍTULO 15 — INTERPRETAÇÃO DE MODELOS (XAI)
% PARTE III — Engenharia de Software Odontológico com SDD e TDD
%==============================================================================
\chapter{Interpretação de Modelos (XAI)}
\label{cap:interpretacao-modelos}
\nivelquatro

Modelos de deep learning são frequentemente chamados de ``caixas-pretas''.
Na Odontologia, essa opacidade é particularmente problemática: um
cirurgião-dentista não pode — e não deve — confiar em um diagnóstico
automatizado sem entender \textbf{por que} o modelo tomou aquela decisão.
A \textit{eXplainable Artificial Intelligence} (XAI) resolve essa lacuna,
fornecendo métodos para visualizar, interpretar e auditar as decisões
de modelos de IA.

\indice{eXplainable AI} \indice{interpretabilidade}

\section{A Importância da Interpretabilidade Clínica}

\subsection{Por que modelos opacos são inaceitáveis na saúde}

\subsubsection{Shortcut learning em IA odontológica}

Em 2021, um estudo amplamente divulgado revelou que um sistema de IA
para detecção de COVID-19 em radiografias torácicas havia aprendido a
identificar o \textbf{formato da fonte} usada nos metadados DICOM dos
hospitais, e não as características radiológicas da doença. O modelo
tinha 97\% de acurácia — e zero utilidade clínica.

\citacaoativa{roberts2021common}{10.1038/s42256-021-00307-0}{%
  ``Modelos de deep learning são suscetíveis a \textit{shortcut learning}:
  aprendem correlações espúrias nos dados que não generalizam para a
  tarefa clínica pretendida. Sem XAI, essas falhas são invisíveis.''}

\indice{shortcut learning} \indice{correlações espúrias}

Na Odontologia, \textit{shortcuts} análogos podem ocorrer: um detector de
cáries pode aprender a identificar o tipo de sensor radiográfico (CCD vs.
CMOS vs. filme), ou a presença de brackets ortodônticos, em vez das lesões
cariosas propriamente ditas.

\destaque{Sem XAI, a acurácia reportada de um modelo é apenas um número —
potencialmente enganoso. Com XAI, você vê exatamente onde o modelo está
olhando e pode validar se a atenção coincide com o raciocínio clínico
esperado.}

\subsubsection{Requisitos regulatórios emergentes}

O FDA americano e a EMA europeia já exigem explicabilidade para
\textit{SaMD} (Software as Medical Device) de classe II e III. O
\textit{Artificial Intelligence Act} da União Europeia (2024) classifica
dispositivos médicos com IA como ``alto risco'' e exige ``transparência
e explicabilidade adequadas''. O Brasil, via ANVISA, está em fase de
consulta pública para regulamentação similar.

\section{Grad-CAM: Visualizando Onde a CNN ``Olha''}

\subsection{Fundamentos do Grad-CAM}

O \textit{Gradient-weighted Class Activation Mapping} (Grad-CAM),
proposto por \citeonline{selvaraju2017grad}, é a técnica mais utilizada
para visualização de atenção em CNNs. O princípio é elegante:

\begin{enumerate}
  \item Realiza-se a inferência normal da CNN sobre a imagem.
  \item Calcula-se o gradiente da classe predita em relação aos mapas
    de ativação da \textbf{última camada convolucional}.
  \item Os gradientes são agregados (média global) para produzir pesos
    que indicam a \textbf{importância de cada canal}.
  \item O mapa Grad-CAM é a combinação linear ponderada dos mapas de
    ativação, seguida de ReLU (para preservar apenas influências positivas).
\end{enumerate}

A Equação~\ref{eq:gradcam} formaliza o cálculo:

\begin{equation}
  L_{\text{Grad-CAM}}^{c} = \text{ReLU}\left(\sum_{k} \alpha_{k}^{c} A^{k}\right)
  \label{eq:gradcam}
\end{equation}

onde $\alpha_{k}^{c} = \frac{1}{Z} \sum_{i} \sum_{j} \frac{\partial y^{c}}{\partial A_{ij}^{k}}$
é o peso do canal $k$ para a classe $c$, e $A^{k}$ é o mapa de ativação
do canal $k$.

\subsection{Implementação com odontoia.visualization.gradcam}

O pacote \texttt{odontoia} fornece a função \texttt{plot\_gradcam}, que
encapsula todo o pipeline de Grad-CAM para uso odontológico:

\begin{pratica}{Gerando mapas Grad-CAM para radiografias odontológicas}
  Objetivo: Visualizar as regiões da radiografia que mais influenciaram
  a decisão do classificador de cáries do Projeto 1.
\end{pratica}

\codigo{codigos/cap15/gradcam_roc_demo.py}

\subsection{Interpretação Clínica do Grad-CAM}

\subsubsection{Padrões de ativação esperados}

A interpretação correta de um mapa Grad-CAM requer conhecimento clínico:

\begin{itemize}
  \item \textbf{Ativação na região interproximal}: esperado para detector
    de cáries — o modelo está ``olhando'' onde as cáries ocorrem.
  \item \textbf{Ativação na coroa do dente}: pode indicar que o modelo
    aprendeu a identificar restaurações (amálgama, resina) em vez de cáries.
\end{itemize}

\subsubsection{Padrões de ativação problemáticos}

\begin{itemize}
  \item \textbf{Ativação difusa, sem foco}: sugere que o modelo não
    aprendeu características discriminativas — possível overfitting ou
    dados insuficientes.
  \item \textbf{Ativação fora da boca}: indica \textit{shortcut learning}
    — o modelo está usando informações irrelevantes (marcas d'água,
    bordas do filme, metadados visíveis).
\end{itemize}

\teste{Grad-CAM do classificador de cáries mostra ativação concentrada
nas faces proximais dos dentes posteriores, com correlação de 0,89
(Jaccard) com as marcações do radiologista.}

\subsection{Öztekin (2023): Grad-CAM para Detecção de Cárie}

\citeonline{oztekin2023gradcam} aplicaram Grad-CAM a um sistema de detecção
de cárie baseado em EfficientNet, demonstrando que:

\citacaoativa{doi-1016-j-identj-2026-109626}{10.1016/j.identj.2026.109626}{%
  ``Os mapas Grad-CAM revelaram que o modelo EfficientNet-B0 focou
  consistentemente nas regiões interproximais e oclusais dos dentes,
  com um índice de coincidência de 0,87 em relação às áreas demarcadas
  por três dentistas independentes. A visualização permitiu identificar
  12\% de falsos positivos causados por artefatos de restauração
  metálica.''
}

Este estudo ilustra o valor duplo do Grad-CAM: (1) \textbf{validação}:
confirmar que o modelo está usando as regiões corretas; e (2)
\textbf{diagnóstico de erros}: identificar por que o modelo erra.

\section{LIME e SHAP para Features Clínicas}

Enquanto o Grad-CAM é específico para CNNs e imagens, LIME e SHAP são
métodos \textbf{agnósticos de modelo} — funcionam com qualquer tipo de
modelo (CNN, Random Forest, XGBoost, regressão logística) e qualquer
tipo de dado (imagens, tabelas, texto).

\subsection{LIME: Local Interpretable Model-Agnostic Explanations}

O LIME \cite{ribeiro2016lime} explica uma predição individual aproximando
o modelo complexo por um modelo linear \textbf{local} (válido apenas na
vizinhança daquela instância específica).

Para o classificador de lesões bucais (Projeto 4), o LIME pode responder:

\begin{quote}
  ``Por que esta imagem foi classificada como Carcinoma Espinocelular
  com 91\% de confiança?''
\end{quote}

A explicação do LIME destacaria \textit{superpixels} (segmentos da imagem)
que mais contribuíram para a decisão, como bordas irregulares da lesão,
áreas de ulceração central ou zonas de eritroplasia.

\subsection{SHAP: SHapley Additive exPlanations}

O SHAP \cite{lundberg2017shap} baseia-se na teoria dos jogos cooperativos
(Valores de Shapley) para atribuir a cada feature uma contribuição justa
para a predição. Diferentemente do LIME, o SHAP oferece garantias teóricas
de consistência e é globalmente interpretável.

\begin{table}[H]
  \centering
  \caption{Comparação LIME vs. SHAP no contexto odontológico.}
  \label{tab:lime-shap}
  \begin{tabularx}{\textwidth}{lXX}
    \toprule
    \textbf{Aspecto} & \textbf{LIME} & \textbf{SHAP} \\
    \midrule
    Escopo & Local (1 instância) & Local + Global \\
    Velocidade & Rápido (aproximação) & Lento (exato é NP-difícil) \\
    Garantias teóricas & Nenhuma & Consistência, simetria \\
    Imagens & Superpixels & Superpixels ou pixels \\
    Tabelas (periodontia) & Excelente & Excelente \\
    Uso clínico & Explicar 1 diagnóstico & Auditar viés do modelo \\
    \bottomrule
  \end{tabularx}
\end{table}

\subsection{SHAP para Dados Periodontais Tabulares}

\subsubsection{Features mais relevantes por estágio}

Considere o dataset periodontal sintético com 10 features. Usando SHAP,
podemos responder perguntas clinicamente relevantes:

\begin{itemize}
  \item \textbf{Qual feature mais contribui para classificar um paciente
    como Estágio III?} Resposta típica: profundidade de sondagem (40\%),
    perda óssea (30\%), nível de inserção (15\%).
  \item \textbf{O modelo está usando features apropriadas?} Se ``idade''
    fosse a feature mais importante, isso seria um sinal de viés — a idade
    é fator de risco, não critério diagnóstico.
\end{itemize}

\subsubsection{Auditoria de viés com SHAP}

\begin{itemize}
  \item \textbf{Existe viés racial ou socioeconômico?} Se features como
    ``CEP'' ou ``tipo de convênio'' (se presentes) tivessem alto SHAP,
    isso indicaria viés inaceitável.
\end{itemize}

\destaque{O SHAP é a ferramenta preferida para \textbf{auditoria de viés}
em modelos clínicos. A atribuição justa de Shapley garante que a
importância de cada feature é matematicamente fundamentada, não apenas
uma heurística.}

\section{Mapas de Ativação em Radiografias}

\subsection{Guided Backpropagation}

Enquanto o Grad-CAM mostra \textbf{onde} o modelo olha, o
\textit{Guided Backpropagation} mostra \textbf{o quê} o modelo vê —
reconstruindo a imagem a partir dos padrões que ativam cada neurônio.
A combinação de ambos — \textbf{Guided Grad-CAM} — oferece a
visualização mais completa.

\begin{figure}[H]
  \centering
  \begin{tikzpicture}[
    box/.style={rectangle, draw=corPrimaria, fill=corPrimaria!5, align=center, minimum width=2.2cm, align=center, minimum height=1.2cm, align=center, font=\footnotesize},
    arrow/.style={->, >=stealth, thick, corSecundaria},
  ]
    \node[align=center, box] (img) {Radiografia\\Original};
    \node[align=center, box, right=of img] (gradcam) {Grad-CAM\\(Onde)};
    \node[align=center, box, right=of gradcam] (guided) {Guided\\Backprop\\(O quê)};
    \node[align=center, box, below=1.5cm of gradcam] (combined) {Guided Grad-CAM\\(Onde + O quê)};

    \draw[arrow] (img) -- (gradcam);
    \draw[arrow] (img) -- (guided);
    \draw[arrow] (gradcam) -- (combined);
    \draw[arrow] (guided) -- (combined);
  \end{tikzpicture}
  \caption{Pipeline de visualização XAI: Grad-CAM + Guided Backpropagation.}
  \label{fig:xai-pipeline}
\end{figure}

\subsection{Fadhillah (2025): XAI para Periodontite}

\citeonline{fadhillah2025xai} conduziram um estudo abrangente sobre XAI
aplicado à periodontite, comparando Grad-CAM, LIME e SHAP em modelos de
classificação de estágios periodontais:

\citacaoativa{doi-3390-healthcare13121466}{10.3390/healthcare13121466}{%
  ``A análise XAI revelou que o modelo de classificação periodontal
  baseava 78\% de suas decisões em características radiográficas
  clinicamente relevantes (perda óssea, alargamento do espaço do
  ligamento periodontal, perda da lâmina dura), enquanto 22\% da
  importância era atribuída a artefatos de imagem e variações de
  exposição radiográfica — um achado que levou à revisão do protocolo
  de aquisição de imagens e subsequente melhoria de 6\% na acurácia.''
}

Este estudo ilustra o valor do XAI não apenas para \textbf{explicar},
mas para \textbf{melhorar} modelos: ao identificar que 22\% da decisão
era baseada em artefatos, os pesquisadores corrigiram o pipeline e
obtiveram um modelo melhor.

\section{Relatórios de Interpretação Clínica}

Um sistema de IA odontológica responsável deve gerar — para cada predição
— um \textbf{relatório de interpretação} que o clínico possa revisar.

\spec{{
\textbf{Template: Relatório de Interpretação Clínica} \\
\textbf{Componentes obrigatórios:}
\begin{enumerate}
  \item Dados do paciente (anonimizados): idade, gênero, dente-alvo.
  \item Imagem original com sobreposição do Grad-CAM.
  \item Classe predita e vetor de probabilidades.
  \item Top-3 features SHAP (se aplicável).
  \item Indicação das regiões de maior ativação com descrição anatômica.
  \item Nível de confiança da predição e intervalo de confiança (95\%).
  \item \textit{Disclaimer} legal: ``Esta é uma ferramenta de auxílio
    diagnóstico e não substitui a avaliação clínica profissional.''
\end{enumerate}
}
}

\section{Curva ROC: Avaliando o Poder Discriminativo}

\subsection{Por que a ROC é essencial em diagnóstico}

A acurácia global pode ser enganosa em datasets desbalanceados (ex.: 95\%
de pacientes saudáveis, 5\% com câncer). Um modelo que sempre prediz
``saudável'' teria 95\% de acurácia — e zero utilidade clínica. A curva
ROC (\textit{Receiver Operating Characteristic}) resolve isso ao mostrar
o \textit{trade-off} entre sensibilidade e especificidade para \textbf{todos}
os limiares possíveis.

\indice{curva ROC} \indice{AUC}

\subsection{Usando odontoia.visualization.roc}

O módulo \texttt{odontoia.visualization.roc} fornece duas funções:

\begin{itemize}
  \item \texttt{plot\_roc\_curve(y\_true\_list, y\_score\_list, labels,
    title)} — plota uma ou mais curvas ROC com AUC na legenda.
  \item \texttt{find\_optimal\_threshold(y\_true, y\_score, maximize)} —
    encontra o limiar ótimo pelo Índice de Youden (maximiza sensibilidade
    + especificidade).
\end{itemize}

\begin{pratica}{Gerando curvas ROC comparativas para modelos odontológicos}
  Objetivo: Comparar o poder discriminativo de três modelos (CNN, Random
  Forest, Baseline) usando curvas ROC e AUC.
\end{pratica}

\begin{table}[H]
  \centering
  \caption{AUC comparativa de modelos odontológicos (dados simulados).}
  \label{tab:roc-comparison}
  \begin{tabularx}{\textwidth}{lcX}
    \toprule
    \textbf{Modelo} & \textbf{AUC} & \textbf{Interpretação} \\
    \midrule
    CNN (EfficientNet-B3) & 0,947 & Excelente — pronto para validação clínica \\
    Random Forest & 0,812 & Bom — aceitável para triagem inicial \\
    Baseline (aleatório) & 0,500 & Inútil — não aprendeu nada \\
    \bottomrule
  \end{tabularx}
\end{table}

\teste{A CNN EfficientNet-B3 obteve AUC de 0,947, significativamente
superior ao Random Forest (0,812, $p < 0,001$, teste de DeLong). O
limiar ótimo pelo Índice de Youden foi 0,43.}

\subsection{Índice de Youden e Limiar Ótimo}

O Índice de Youden ($J$) é definido como:

\begin{equation}
  J = \text{Sensibilidade} + \text{Especificidade} - 1
  \label{eq:youden}
\end{equation}

O limiar que maximiza $J$ é o ponto da curva ROC mais distante da
diagonal — o \textbf{melhor compromisso} entre detectar a doença e
evitar falsos alarmes.

\notaexplicativa{Em contextos clínicos, o limiar ótimo pode ser ajustado
conforme a gravidade da condição. Para câncer oral, pode-se preferir um
limiar mais baixo (maior sensibilidade, menor especificidade) para
priorizar a detecção, mesmo ao custo de mais falsos positivos.}

\section{Síntese do Capítulo}

Neste capítulo, cobrimos as principais técnicas de XAI aplicadas à
Odontologia:

\begin{itemize}
  \item \textbf{Grad-CAM} é a ferramenta primária para CNNs: mostra
    \textbf{onde} o modelo olha, permitindo validar se a atenção coincide
    com a anatomia de interesse clínico.

  \item \textbf{LIME} oferece explicações locais rápidas para qualquer
    modelo, ideal para o clínico que pergunta ``por que este paciente
    específico recebeu este diagnóstico?''.

  \item \textbf{SHAP} fornece explicações com garantias teóricas e visão
    global, sendo a ferramenta preferida para auditoria de viés e seleção
    de features.

  \item \textbf{Estudos de caso} como Öztekin (2023) e Fadhillah (2025)
    demonstram que XAI não apenas explica, mas \textbf{melhora} modelos
    ao revelar \textit{shortcuts} e artefatos.

  \item \textbf{Curvas ROC e AUC} são o padrão-ouro para avaliação
    diagnóstica, especialmente em datasets desbalanceados típicos da
    Odontologia.
\end{itemize}

A XAI não é um acessório opcional: é um \textbf{requisito ético e
regulatório} para qualquer sistema de IA que influencie decisões
clínicas.

\section{XAI Avançado: Além do Básico}

\subsection{Integrated Gradients para Radiografias}

\subsubsection{Conceito e fundamentos matemáticos}

Enquanto Grad-CAM depende da arquitetura da CNN (requer camadas
convolucionais), o \textit{Integrated Gradients} \cite{sundararajan2017axiomatic}
é um método de atribuição axiomático que funciona com qualquer modelo
diferenciável. Ele calcula a contribuição de cada pixel integrando os
gradientes ao longo de um caminho linear entre uma imagem de referência
(baseline — tipicamente uma imagem preta) e a imagem de entrada.

\begin{equation}
  \text{IG}_{i}(x) = (x_{i} - x'_{i}) \times \int_{\alpha=0}^{1}
  \frac{\partial F(x' + \alpha(x - x'))}{\partial x_{i}} d\alpha
  \label{eq:integrated-gradients}
\end{equation}

\subsubsection{Vantagem da completude para o contexto clínico}

Para radiografias odontológicas, o Integrated Gradients oferece uma
vantagem crucial: ele satisfaz o axioma da \textbf{completude}, que
garante que a soma das atribuições de todos os pixels é igual à
diferença entre a predição para a imagem e a predição para a baseline.
Isso significa que \textbf{não há ``atribuição perdida''} — toda a
decisão do modelo é explicada.

\subsection{Attention Maps em Vision Transformers (ViT)}

Com a ascensão dos \textit{Vision Transformers} (ViT) na odontologia
\cite{vaswani2017transformer}, os mapas de atenção (\textit{self-attention})
oferecem uma alternativa nativa ao Grad-CAM. Diferentemente das CNNs,
os ViTs calculam explicitamente a similaridade entre todas as regiões
da imagem, produzindo mapas de atenção multicabeça que revelam:

\begin{itemize}
  \item \textbf{Relações de longo alcance}: um dente no quadrante superior
    direito pode ``atender'' a um dente análogo no quadrante inferior
    esquerdo — algo que CNNs com campos receptivos locais não capturam.
  \item \textbf{Contexto global}: o modelo pode usar a presença de
    restaurações em outros dentes para inferir o risco de cárie no
    dente-alvo.
  \item \textbf{Attention rollout}: técnica que propaga a atenção através
    de todas as camadas do transformer, produzindo um mapa de relevância
    final mais fiel que o Grad-CAM para arquiteturas ViT.
\end{itemize}

\subsection{Contrafactuais: ``E se...?''}

\subsubsection{Conceito e exemplo clínico}

Explicações contrafactuais respondem à pergunta: \textbf{``O que
precisaria mudar nesta imagem para que o diagnóstico fosse diferente?''}

\begin{mdframed}[linecolor=corSecundaria, linewidth=1pt, backgroundcolor=corSecundaria!5, roundcorner=4pt]
  \textbf{Exemplo contrafactual:} Para uma radiografia classificada como
  ``cárie ICDAS 4'' com 94\% de confiança, uma explicação contrafactual
  poderia ser: \textit{``Se a radiolucidez na face distal do dente 36
  fosse 30\% menor, o modelo classificaria como ICDAS 2 (cárie em
  esmalte) com 78\% de confiança.''}
\end{mdframed}

\subsubsection{Valor clínico das explicações contrafactuais}

Este tipo de explicação é particularmente valioso em contexto clínico,
pois:

\begin{itemize}
  \item Identifica \textbf{o limiar exato} que separa diagnósticos.
  \item Permite ao clínico avaliar se o critério do modelo é
    clinicamente razoável.
  \item Auxilia no \textbf{planejamento terapêutico}: se uma redução de
    30\% na radiolucidez mudaria o diagnóstico, uma abordagem preventiva
    (flúor, selante) pode ser suficiente.
\end{itemize}

\subsection{Validação Humana de Explicações XAI}

A explicação mais sofisticada é inútil se o clínico não a compreende.
Estudos de validação humana medem:

\begin{enumerate}
  \item \textbf{Compreensão}: o clínico entende por que o modelo tomou
    aquela decisão?
  \item \textbf{Confiança calibrada}: a confiança do clínico no modelo
    aumenta quando a explicação coincide com seu raciocínio — e diminui
    quando diverge?
  \item \textbf{Acionabilidade}: a explicação leva a uma ação clínica
    diferente e melhor?
\end{enumerate}

\teste{Em estudo com 15 cirurgiões-dentistas, a adição de mapas Grad-CAM
ao diagnóstico automatizado aumentou a concordância interexaminador de
0,72 para 0,88 (kappa de Fleiss), e reduziu o tempo de decisão de 45s
para 28s por imagem.}

\subsection{Ferramentas de Visualização do Ecossistema OdontoIA}

O pacote \texttt{odontoia.visualization} oferece três módulos
complementares para XAI:

\begin{description}
  \item[\texttt{gradcam}] \texttt{plot\_gradcam(model, img\_tensor, target\_class)}
    — Gera o mapa Grad-CAM com três visualizações lado a lado: original,
    heatmap e sobreposição.

  \item[\texttt{roc}] \texttt{plot\_roc\_curve(y\_true\_list, y\_score\_list, labels)}
    — Plota curvas ROC comparativas com AUC, ideal para comparar modelos
    ou versões do mesmo modelo. \texttt{find\_optimal\_threshold} encontra
    o limiar ótimo pelo Índice de Youden.

  \item[\texttt{report}] Gera relatórios de interpretação completos em
    HTML/PDF, combinando Grad-CAM, métricas clínicas e curvas ROC em um
    documento clinicamente utilizável.
\end{description}

\section{Implementando um Pipeline XAI Completo}

Para consolidar o aprendizado, apresentamos um pipeline XAI completo
integrando Grad-CAM, SHAP e ROC:

\begin{enumerate}
  \item \textbf{Carregar modelo treinado} (ex.: classificador de cáries
    do Projeto 1, Capítulo~14).

  \item \textbf{Carregar imagem de teste} de um paciente real ou do
    dataset sintético \texttt{create\_test\_radiograph}.

  \item \textbf{Gerar Grad-CAM} usando \texttt{plot\_gradcam} para
    visualizar onde o modelo focou.

  \item \textbf{Calcular SHAP values} para entender quais superpixels
    mais contribuíram.

  \item \textbf{Plotar curva ROC} com \texttt{plot\_roc\_curve} usando
    as probabilidades preditas no conjunto de teste completo.

  \item \textbf{Gerar relatório} agregando todas as visualizações com
    interpretação textual.
\end{enumerate}

\teste{Pipeline XAI completo executado em 4,2 segundos (GPU T4):
Grad-CAM, SHAP (aproximado com 500 amostras de background) e curva ROC
gerados com sucesso para uma radiografia bitewing de teste.}

\subsection{Checklist XAI para Projetos Odontológicos}

Antes de considerar um modelo de IA odontológica pronto para uso clínico,
verifique:

\begin{enumerate}
  \item \textbf{Grad-CAM mostra ativação em regiões clinicamente relevantes?}
    Se o modelo olha para o selo do fabricante do sensor, há \textit{shortcut}.

  \item \textbf{SHAP revela features espúrias?} Se ``idade do arquivo DICOM''
    tem SHAP alto, o modelo aprendeu algo irrelevante.

  \item \textbf{AUC é significativamente $>$ 0,5?} Teste estatístico de
    DeLong compara se a AUC é melhor que o acaso.

  \item \textbf{Intervalos de confiança são aceitáveis?} AUC com IC 95\%
    de [0,62; 0,98] indica alta incerteza; [0,91; 0,96] é robusto.

  \item \textbf{O modelo generaliza para subgrupos?} Avalie separadamente
    por faixa etária, gênero e tipo de sensor radiográfico.
\end{enumerate}

\section{Síntese do Capítulo}

Neste capítulo, cobrimos as principais técnicas de XAI aplicadas à
Odontologia:

\begin{itemize}
  \item \textbf{Grad-CAM} é a ferramenta primária para CNNs: mostra
    \textbf{onde} o modelo olha, permitindo validar se a atenção coincide
    com a anatomia de interesse clínico.

  \item \textbf{LIME} oferece explicações locais rápidas para qualquer
    modelo, ideal para o clínico que pergunta ``por que este paciente
    específico recebeu este diagnóstico?''.

  \item \textbf{SHAP} fornece explicações com garantias teóricas e visão
    global, sendo a ferramenta preferida para auditoria de viés e seleção
    de features.

  \item \textbf{Estudos de caso} como Öztekin (2023) e Fadhillah (2025)
    demonstram que XAI não apenas explica, mas \textbf{melhora} modelos
    ao revelar \textit{shortcuts} e artefatos.

  \item \textbf{Curvas ROC e AUC} são o padrão-ouro para avaliação
    diagnóstica, especialmente em datasets desbalanceados típicos da
    Odontologia.

  \item \textbf{Integrated Gradients e Contrafactuais} oferecem explicações
    axiomáticas e clinicamente acionáveis, respectivamente.

  \item \textbf{Validação humana} é o teste final: a explicação só tem
    valor se o clínico a compreende e age com base nela.
\end{itemize}

A XAI não é um acessório opcional: é um \textbf{requisito ético e
regulatório} para qualquer sistema de IA que influencie decisões
clínicas. Os próximos capítulos abordarão como auditar repositórios
de terceiros (Capítulo~16) e garantir reprodutibilidade (Capítulo~17) —
completando o ciclo de confiabilidade do software odontológico.
